\documentclass{article}
\usepackage[T1]{fontenc}
\usepackage{lmodern}
\usepackage{graphicx}
\usepackage{amsmath,amssymb}
\usepackage[a4paper,margin=2cm]{geometry}
\usepackage[absolute,overlay]{textpos}
\setlength{\TPHorizModule}{1mm}
\setlength{\TPVertModule}{1mm}
\usepackage[indonesian]{babel}
\usepackage{hyperref}
\setlength{\emergencystretch}{2em}
% unit: Halaman 3

\begin{document}

% === Halaman 3 (llm) ===

\begin{itemize}
    \item Apakah \textit{unbreakable cipher} memang benar-benar ada?
    Jawaban: ada!
    \item Apa syarat sebuah \textit{cipher} disebut \textit{unbreakable cipher}?
    Jawaban:
    \begin{enumerate}
        \item Kunci harus acak sejati (\textit{truly random}).
        \item Panjang kunci = panjang plainteks
        \item Kunci hanya boleh digunakan sekali, tidak boleh digunakan ulang untuk mengenkripsi pesan yang lain
    \end{enumerate}
    \item Acak sejati artinya tidak dapat diprediksi nilainya dan tidak dapat diulang kembali pembangkitannya. Ini berbeda dengan \textit{pseudo random} yang bisa diulang kembali pembangkitan nilai acak asalkan umpan (\textit{seed}) nya sama.
    \item Akibat 1 dan 2: plainteks yang sama tidak selalu menghasilkan cipherteks yang sama
\end{itemize}

\end{document}
